\begin{textsample}{POS Dim 3 – rj – Score 19.00 – rj/rj\_inf\_alfabetizar\_2.txt}  \label{ex:f3_pos_013}
Nas sociedades letradas, desde os primeiros \textbf{meses}, as \textbf{crianças} estão em permanente contato com a linguagem escrita.
É por meio desse contato diversificado em seu ambiente social que descobrem o aspecto funcional da comunicação escrita, desenvolvendo interesse e \textbf{curiosidade} por essa linguagem.
Diante do ambiente de letramento em que vivem, podem fazer, a partir de dois ou três anos de idade, uma série de perguntas, como “ O que está escrito aqui? ”, ou “ O que isto quer dizer? ”, indicando sua reflexão sobre a função e o significado da escrita, ao perceberem que ela representa algo.
Sabe -se que, para aprender a escrever, a \textbf{criança} terá de lidar com dois processos de aprendizagem paralelos : o da natureza do sistema de escrita da língua – o que a escrita representa e como – e o das características da linguagem que se usa para escrever.
A aprendizagem da linguagem escrita está intrinsecamente associada ao contato com textos diversos, para que as \textbf{crianças} possam construir sua capacidade de ler, e às práticas de escrita, para que possam desenvolver a capacidade de escrever autonomamente.
A observação e a análise das produções escritas das \textbf{crianças} revelam que elas tomam consciência, gradativamente, das características formais dessa linguagem.
Constata -se que, desde muito \textbf{pequenas}, as \textbf{crianças} podem usar o \textbf{lápis} e o papel para imprimir \textbf{marcas}, \textbf{imitando} a escrita dos mais velhos, assim como se utilizam de \textbf{livros} de literatura, revistas, jornais, gibis, rótulos etc. para “ ler ” o que está escrito.
Não é raro observar \textbf{crianças} muito \textbf{pequenas}, que têm contato com material escrito, folheando um \textbf{livro}, \textbf{emitindo} \textbf{sons} e fazendo \textbf{gestos}, como se estivessem lendo.
A escrita da \textbf{criança} não resulta de simples \textbf{cópia} de um modelo externo, mas é um processo de construção pessoal, como diz Emilia Ferreiro ( 2000 ).
Essa autora também nos \textbf{convida} a defender a importância de a Educação \textbf{Infantil} assumir o compromisso com um ambiente alfabetizador inundado da linguagem escrita em nossas salas das \textbf{creches} e \textbf{pré-escolas}.
Nesse ambiente, devem constar os textos orais das \textbf{crianças} registrados pelo professor na forma escrita, além de haver a presença constante e pujante de bilhetes, cartas, receitas, \textbf{livros} de literatura \textbf{infantil}, entre outros.
Nesse sentido, a linguagem escrita aparece no cotidiano do trabalho como uma das múltiplas linguagens ( oral, \textbf{mímica}, desenho, colagem, \textbf{pintura} etc. ) que as \textbf{crianças} utilizam para livremente exprimir suas vivências, experiências, sentimentos, \textbf{desejos} e necessidades.

% matched lemmas: convidar, creche, criança, curiosidade, cópia, desejo, emitir, gesto, imitar, infantil, livro, lápis, marca, mês, mímico, pequeno, pintura, pré-escola, som
\end{textsample}
